\documentclass[12pt, a4paper]{article}

\usepackage[utf8]{inputenc}
\usepackage[T1]{fontenc}
\usepackage[polish]{babel}
\usepackage{geometry}
\usepackage{graphicx}
\usepackage{hyperref}
\usepackage{listings}
\usepackage{xcolor}
\usepackage{tikz}
\usepackage{float}
\usepackage{enumitem}
\usepackage{booktabs}
\usepackage{tabularx}
\usepackage{fancyhdr}
\usepackage{parskip}

\usetikzlibrary{shapes, shapes.geometric, arrows.meta, positioning, fit, backgrounds, calc}

\geometry{a4paper, top=2.5cm, bottom=2.5cm, left=2.5cm, right=2.5cm}

\pagestyle{fancy}
\fancyhf{}
\rhead{System Rezerwacji Terminów Lekarskich}
\lhead{ZTI -- Projekt}
\rfoot{\thepage}

\definecolor{codebg}{RGB}{245,245,245}
\definecolor{codegreen}{rgb}{0,0.5,0}
\definecolor{codegray}{rgb}{0.5,0.5,0.5}
\definecolor{codeblue}{rgb}{0,0,0.7}
\definecolor{codepurple}{rgb}{0.5,0,0.5}

\lstdefinestyle{java}{
  backgroundcolor=\color{codebg},
  commentstyle=\color{codegreen},
  keywordstyle=\color{codeblue}\bfseries,
  stringstyle=\color{codepurple},
  basicstyle=\ttfamily\footnotesize,
  breaklines=true,
  keepspaces=true,
  numbers=left,
  numberstyle=\tiny\color{codegray},
  showstringspaces=false,
  tabsize=2,
  language=Java,
  frame=single,
  rulecolor=\color{gray!30}
}

\lstdefinestyle{bash}{
  backgroundcolor=\color{codebg},
  basicstyle=\ttfamily\footnotesize,
  breaklines=true,
  frame=single,
  rulecolor=\color{gray!30},
  numbers=none
}

\hypersetup{colorlinks=true, linkcolor=blue, urlcolor=blue}

% --- tikz styles ---
\tikzset{
  classbox/.style={rectangle, draw=black, thick, fill=blue!8,
    minimum width=4cm, minimum height=0.7cm, align=center, font=\small},
  classattr/.style={rectangle, draw=black, fill=white,
    minimum width=4cm, align=left, font=\footnotesize, inner sep=4pt},
  actor/.style={font=\small, align=center},
  usecase/.style={ellipse, draw, fill=yellow!20, align=center, font=\footnotesize,
    minimum width=2.8cm, minimum height=0.8cm},
  seqbox/.style={rectangle, draw, fill=blue!15, minimum width=2cm,
    minimum height=0.6cm, align=center, font=\footnotesize},
  lifeline/.style={dashed, gray},
  msgarr/.style={-Stealth, thick},
  retarr/.style={-Stealth, dashed}
}

% ============================================================
\begin{document}

\begin{titlepage}
  \centering
  \vspace*{2cm}
  {\LARGE\bfseries Zaawansowane Technologie Internetowe\par}
  \vspace{0.5cm}
  {\large Dokumentacja Projektu\par}
  \vspace{2cm}
  {\huge\bfseries System Rezerwacji Terminów Lekarskich\par}
  \vspace{1cm}
  {\large \textit{DocBooking}\par}
  \vspace{3cm}
  \begin{tabular}{ll}
    \textbf{Autor:}       & Artur Saganowski \\[4pt]
    \textbf{Technologie:} & Spring Boot 3.4, Angular 20, PostgreSQL \\[4pt]
    \textbf{Data:}        & \today \\
  \end{tabular}
  \vfill
\end{titlepage}

\tableofcontents
\newpage

% ============================================================
\section{Opis funkcjonalności aplikacji}

\subsection{Cel systemu}

\textbf{DocBooking} to aplikacja webowa klasy klient-serwer umożliwiająca pacjentom rezerwację wizyt lekarskich online. System zapewnia podgląd dostępnych terminów w czasie rzeczywistym dzięki technologii WebSocket, eliminując problem podwójnej rezerwacji tego samego slotu.

\subsection{Funkcjonalności części serwerowej}

\begin{enumerate}[label=\arabic*.]
  \item \textbf{Zarządzanie użytkownikami} -- rejestracja pacjentów, logowanie, generowanie i walidacja tokenów JWT.
  \item \textbf{Zarządzanie lekarzami} -- pobieranie listy lekarzy, filtrowanie po specjalizacji.
  \item \textbf{Zarządzanie slotami} -- tworzenie slotów przez lekarza, pobieranie dostępnych terminów, zmiana statusu slotu (\texttt{AVAILABLE} / \texttt{BEING\_RESERVED} / \texttt{BOOKED}).
  \item \textbf{Rezerwacja wizyt} -- tworzenie rezerwacji, anulowanie wizyty, historia wizyt pacjenta.
  \item \textbf{WebSocket} -- rozgłaszanie zmian statusu slotów do wszystkich klientów przeglądających terminy tego samego lekarza.
  \item \textbf{Spring AOP} -- automatyczne logowanie każdego wywołania serwisu (czas wykonania, użytkownik, błędy) oraz audit zdarzeń rezerwacji.
  \item \textbf{Inicjalizacja danych} -- seedowanie bazy przy pierwszym uruchomieniu: 3 lekarzy z terminami na 5 dni roboczych.
\end{enumerate}

\subsection{Funkcjonalności klienta (Angular 20)}

\begin{enumerate}[label=\arabic*.]
  \item \textbf{Rejestracja i logowanie} -- formularze z walidacją, token JWT w \texttt{localStorage}.
  \item \textbf{Lista lekarzy} -- karty Angular Material z danymi lekarza.
  \item \textbf{Widok terminów} -- siatka slotów z live aktualizacją przez WebSocket.
  \item \textbf{Rezerwacja} -- dwuetapowy proces: oznaczenie slotu jako „rezerwowany", następnie zapis wizyty.
  \item \textbf{Historia wizyt} -- lista wizyt z możliwością anulowania.
  \item \textbf{Guard i Interceptor} -- ochrona tras i automatyczne dołączanie tokenu JWT.
\end{enumerate}

% ============================================================
\section{Architektura systemu}

\subsection{Diagram komponentów}

\begin{figure}[H]
\centering
\begin{tikzpicture}[node distance=3.5cm]
  \node[rectangle, draw, rounded corners, fill=blue!15, minimum width=3cm,
        minimum height=1.2cm, align=center, font=\small] (fe)
        {Angular 20\\Frontend\\:4200};
  \node[rectangle, draw, rounded corners, fill=green!15, minimum width=3cm,
        minimum height=1.2cm, align=center, font=\small, right=of fe] (be)
        {Spring Boot\\Backend\\:8080};
  \node[cylinder, shape border rotate=90, draw, fill=orange!20,
        minimum width=2.5cm, minimum height=1.4cm, align=center,
        font=\small, right=of be] (db) {PostgreSQL\\:5432};

  \draw[-Stealth, thick] (fe) -- node[above, font=\footnotesize]{REST / WS} (be);
  \draw[-Stealth, thick] (be) -- node[above, font=\footnotesize]{JPA} (db);
  \draw[-Stealth, dashed, bend left=40] (be) to
        node[above, font=\footnotesize]{STOMP} (fe);
\end{tikzpicture}
\caption{Diagram komponentów systemu DocBooking}
\end{figure}

\subsection{Stos technologiczny}

\begin{table}[H]
\centering
\begin{tabularx}{\textwidth}{llX}
\toprule
\textbf{Warstwa} & \textbf{Technologia} & \textbf{Rola} \\
\midrule
Frontend  & Angular 20            & SPA, routing, formularze \\
UI        & Angular Material      & Komponenty (karty, przyciski, snackbar) \\
HTTP      & HttpClient + Interceptor & Żądania REST z nagłówkiem JWT \\
WebSocket & STOMP.js + SockJS     & Subskrypcja live aktualizacji \\
Backend   & Spring Boot 3.4.3     & REST API, WebSocket broker \\
Security  & Spring Security + JWT & Autentykacja bezstanowa \\
ORM       & Spring Data JPA       & Mapowanie obiektowo-relacyjne \\
AOP       & Spring AOP            & Logowanie aspektowe \\
Baza      & PostgreSQL 16         & Relacyjna baza danych \\
Infra     & Docker + Compose      & Konteneryzacja wszystkich serwisów \\
\bottomrule
\end{tabularx}
\caption{Stos technologiczny aplikacji}
\end{table}

% ============================================================
\section{Diagramy UML}

\subsection{Diagram klas -- model dziedziny}

\begin{figure}[H]
\centering
\begin{tikzpicture}[node distance=0.5cm]

% --- User ---
\node[classbox] (uhdr) {User};
\node[classattr, below=0 of uhdr, minimum width=4cm] (uattr)
  {\texttt{+id: Long}\\
   \texttt{+username: String}\\
   \texttt{+email: String}\\
   \texttt{-password: String}\\
   \texttt{+role: \{PATIENT,DOCTOR\}}};

% --- Doctor ---
\node[classbox, right=3cm of uhdr] (dhdr) {Doctor};
\node[classattr, below=0 of dhdr, minimum width=4cm] (dattr)
  {\texttt{+id: Long}\\
   \texttt{+firstName: String}\\
   \texttt{+lastName: String}\\
   \texttt{+specialization: String}\\
   \texttt{+phoneNumber: String}};

% --- Patient ---
\node[classbox, below=2.5cm of uhdr] (phdr) {Patient};
\node[classattr, below=0 of phdr, minimum width=4cm] (pattr)
  {\texttt{+id: Long}\\
   \texttt{+firstName: String}\\
   \texttt{+lastName: String}\\
   \texttt{+dateOfBirth: LocalDate}};

% --- TimeSlot ---
\node[classbox, right=3cm of dhdr] (thdr) {TimeSlot};
\node[classattr, below=0 of thdr, minimum width=4cm] (tattr)
  {\texttt{+id: Long}\\
   \texttt{+startTime: LocalDateTime}\\
   \texttt{+endTime: LocalDateTime}\\
   \texttt{+status: SlotStatus}};

% --- Appointment ---
\node[classbox, below=2.5cm of thdr] (ahdr) {Appointment};
\node[classattr, below=0 of ahdr, minimum width=4cm] (aattr)
  {\texttt{+id: Long}\\
   \texttt{+notes: String}\\
   \texttt{+status: AppointmentStatus}\\
   \texttt{+createdAt: LocalDateTime}};

% --- Relacje ---
\draw[-Stealth] (uhdr.east) -- node[above, font=\tiny]{1\quad 0..1} (dhdr.west);
\draw[-Stealth] (uhdr.south) -- node[right, font=\tiny]{1\quad 0..1} (phdr.north);
\draw[-Stealth] (dhdr.east) -- node[above, font=\tiny]{1\quad *} (thdr.west);
\draw[-Stealth] (thdr.south) -- node[right, font=\tiny]{1\quad 0..1} (ahdr.north);
\draw[-Stealth] (phdr.east) -- node[above, font=\tiny]{1\quad *} (ahdr.west);

\end{tikzpicture}
\caption{Diagram klas modelu dziedziny (backend)}
\end{figure}

\subsection{Diagram sekwencji -- rezerwacja wizyty}

\begin{figure}[H]
\centering
\begin{tikzpicture}[x=2.4cm, y=-0.7cm, font=\footnotesize]

% --- Nagłówki uczestników ---
\foreach \x/\lbl in {0/Pacjent, 1/Angular, 2/Backend, 3/WebSocket, 4/PostgreSQL} {
  \node[seqbox] at (\x, 0) {\lbl};
  \draw[lifeline] (\x, 0.4) -- (\x, 10.5);
}

% --- Wiadomości ---
\draw[msgarr] (0,1) -- node[above]{\textit{klik Zarezerwuj}} (1,1);
\draw[msgarr] (1,1.8) -- node[above]{POST /slots/\{id\}/reserve} (2,1.8);
\draw[msgarr] (2,2.6) -- node[above]{UPDATE BEING\_RESERVED} (4,2.6);
\draw[retarr] (4,3.2) -- node[above]{ok} (2,3.2);
\draw[msgarr] (2,4.0) -- node[above]{publish /topic/slots/\{id\}} (3,4.0);
\draw[msgarr] (3,4.8) -- node[above]{slot update (BEING\_RESERVED)} (1,4.8);

\draw[msgarr] (1,5.6) -- node[above]{POST /appointments} (2,5.6);
\draw[msgarr] (2,6.4) -- node[above]{INSERT appointment} (4,6.4);
\draw[msgarr] (2,7.2) -- node[above]{UPDATE BOOKED} (4,7.2);
\draw[retarr] (4,7.8) -- node[above]{ok} (2,7.8);
\draw[msgarr] (2,8.6) -- node[above]{publish /topic/slots/\{id\}} (3,8.6);
\draw[msgarr] (3,9.4) -- node[above]{slot update (BOOKED)} (1,9.4);
\draw[retarr] (2,10.2) -- node[above]{201 Created} (1,10.2);

\end{tikzpicture}
\caption{Diagram sekwencji -- rezerwacja wizyty z live aktualizacją WebSocket}
\end{figure}

\subsection{Diagram przypadków użycia}

\begin{figure}[H]
\centering
\begin{tikzpicture}[node distance=1cm]

% Aktorzy
\node[align=center, font=\small] (pac) at (-4, -3) {
  \begin{tikzpicture}
    \draw[thick] (0,0.4) circle (0.25);
    \draw[thick] (0,0.15) -- (0,-0.4);
    \draw[thick] (-0.3,-0.05) -- (0.3,-0.05);
    \draw[thick] (0,-0.4) -- (-0.2,-0.8);
    \draw[thick] (0,-0.4) -- (0.2,-0.8);
  \end{tikzpicture}\\Pacjent};

\node[align=center, font=\small] (lek) at (4, -3) {
  \begin{tikzpicture}
    \draw[thick] (0,0.4) circle (0.25);
    \draw[thick] (0,0.15) -- (0,-0.4);
    \draw[thick] (-0.3,-0.05) -- (0.3,-0.05);
    \draw[thick] (0,-0.4) -- (-0.2,-0.8);
    \draw[thick] (0,-0.4) -- (0.2,-0.8);
  \end{tikzpicture}\\Lekarz};

% Granica systemu
\draw[thick] (-2.5, 0.5) rectangle (2.5, -6.5);
\node[font=\small\bfseries] at (0, 0.2) {DocBooking};

% Use-case'y
\node[usecase] (reg)  at (0, -0.8)  {Rejestracja};
\node[usecase] (log)  at (0, -1.9)  {Logowanie};
\node[usecase] (list) at (0, -3.0)  {Przeglądanie lekarzy};
\node[usecase] (book) at (0, -4.1)  {Rezerwacja wizyty};
\node[usecase] (canc) at (0, -5.2)  {Anulowanie wizyty};
\node[usecase] (slot) at (0, -6.2)  {Dodawanie terminów};

% Powiązania Pacjent
\draw (pac) -- (reg);
\draw (pac) -- (log);
\draw (pac) -- (list);
\draw (pac) -- (book);
\draw (pac) -- (canc);

% Powiązania Lekarz
\draw (lek) -- (log);
\draw (lek) -- (slot);

\end{tikzpicture}
\caption{Diagram przypadków użycia}
\end{figure}

% ============================================================
\section{Opis kluczowych komponentów}

\subsection{Warstwa bezpieczeństwa -- JWT}

Autentykacja opiera się na tokenach JWT podpisanych kluczem HMAC-SHA256. Po zalogowaniu klient przechowuje token w \texttt{localStorage} i dołącza go do każdego żądania przez interceptor.

\begin{lstlisting}[style=java, caption={Generowanie tokenu JWT -- JwtUtil.java}]
public String generateToken(UserDetails userDetails) {
    return Jwts.builder()
            .subject(userDetails.getUsername())
            .issuedAt(new Date())
            .expiration(new Date(System.currentTimeMillis() + expiration))
            .signWith(getSigningKey())
            .compact();
}
\end{lstlisting}

\subsection{Programowanie aspektowe -- Spring AOP}

Aspekt \texttt{LoggingAspect} automatycznie loguje każde wywołanie warstwy serwisowej bez modyfikacji kodu biznesowego.

\begin{lstlisting}[style=java, caption={Aspekt logowania -- LoggingAspect.java}]
@Aspect
@Component
public class LoggingAspect {

    @Pointcut("execution(* com.zti.doctorreservation.service.*.*(..))")
    public void serviceLayer() {}

    @Around("serviceLayer()")
    public Object logServiceExecution(ProceedingJoinPoint joinPoint)
            throws Throwable {
        String method = joinPoint.getSignature().toShortString();
        long start = System.currentTimeMillis();
        Object result = joinPoint.proceed();
        long elapsed = System.currentTimeMillis() - start;
        log.info("[{}] {} completed in {}ms",
                getCurrentUser(), method, elapsed);
        return result;
    }

    @After("execution(* com.zti.doctorreservation.service"
         + ".AppointmentService.book(..))")
    public void logAppointmentBooked(JoinPoint joinPoint) {
        log.info("[AUDIT] Appointment booked by: {}", getCurrentUser());
    }
}
\end{lstlisting}

\subsection{Komunikacja WebSocket -- STOMP}

\begin{lstlisting}[style=java, caption={Rozgłaszanie aktualizacji slotu -- TimeSlotService.java}]
private void broadcastSlotUpdate(Long doctorId, TimeSlot slot) {
    messagingTemplate.convertAndSend(
        "/topic/slots/" + doctorId, slot);
}
\end{lstlisting}

\begin{lstlisting}[style=java, caption={Subskrypcja WebSocket -- websocket.ts (Angular)}]
subscribeToSlots(doctorId: number): Observable<any> {
    const subject = new Subject<any>();
    this.client.subscribe(
        `/topic/slots/${doctorId}`,
        (msg: IMessage) => subject.next(JSON.parse(msg.body))
    );
    return subject.asObservable();
}
\end{lstlisting}

% ============================================================
\section{REST API}

\begin{table}[H]
\centering
\begin{tabularx}{\textwidth}{llXl}
\toprule
\textbf{Metoda} & \textbf{Ścieżka} & \textbf{Opis} & \textbf{Rola} \\
\midrule
POST & /api/auth/register               & Rejestracja             & Publiczny \\
POST & /api/auth/login                  & Logowanie, zwraca JWT   & Publiczny \\
GET  & /api/doctors                     & Lista lekarzy           & AUTH \\
GET  & /api/doctors/\{id\}              & Dane lekarza            & AUTH \\
GET  & /api/doctors/\{id\}/slots        & Dostępne terminy        & AUTH \\
POST & /api/doctors/\{id\}/slots        & Dodaj termin            & DOCTOR \\
POST & /api/appointments                & Zarezerwuj wizytę       & PATIENT \\
GET  & /api/appointments/my             & Moje wizyty             & PATIENT \\
PUT  & /api/appointments/\{id\}/cancel  & Anuluj wizytę           & PATIENT \\
POST & /api/appointments/slots/\{id\}/reserve & Oznacz slot        & PATIENT \\
\bottomrule
\end{tabularx}
\caption{Endpointy REST API}
\end{table}

% ============================================================
\section{Testy}

\subsection{Testy jednostkowe -- AuthService}

\begin{lstlisting}[style=java, caption={Test rejestracji -- AuthServiceTest.java}]
@ExtendWith(MockitoExtension.class)
class AuthServiceTest {

    @Mock UserRepository userRepository;
    @Mock PatientRepository patientRepository;
    @Mock PasswordEncoder passwordEncoder;
    @Mock JwtUtil jwtUtil;
    @Mock AuthenticationManager authenticationManager;
    @Mock UserDetailsService userDetailsService;
    @InjectMocks AuthService authService;

    @Test
    void register_ShouldCreatePatientAndReturnToken() {
        RegisterRequest request = new RegisterRequest();
        request.setUsername("jan");
        request.setEmail("jan@test.pl");
        request.setPassword("haslo123");
        request.setFirstName("Jan");
        request.setLastName("Nowak");
        request.setRole(User.Role.PATIENT);

        when(userRepository.existsByUsername("jan")).thenReturn(false);
        when(userRepository.existsByEmail("jan@test.pl")).thenReturn(false);
        when(passwordEncoder.encode("haslo123")).thenReturn("$hashed$");
        when(userRepository.save(any())).thenAnswer(i -> {
            User u = i.getArgument(0); u.setId(1L); return u;
        });
        UserDetails ud = mock(UserDetails.class);
        when(userDetailsService.loadUserByUsername("jan")).thenReturn(ud);
        when(jwtUtil.generateToken(ud)).thenReturn("token-abc");

        AuthResponse response = authService.register(request);

        assertEquals("token-abc", response.getToken());
        assertEquals("PATIENT", response.getRole());
        verify(patientRepository).save(any(Patient.class));
    }

    @Test
    void register_ShouldThrow_WhenUsernameExists() {
        RegisterRequest request = new RegisterRequest();
        request.setUsername("jan");
        when(userRepository.existsByUsername("jan")).thenReturn(true);

        assertThrows(IllegalArgumentException.class,
                () -> authService.register(request));
    }
}
\end{lstlisting}

\subsection{Test integracyjny -- kontroler REST}

\begin{lstlisting}[style=java, caption={Test integracyjny -- AuthControllerTest.java}]
@SpringBootTest(webEnvironment = RANDOM_PORT)
@AutoConfigureMockMvc
class AuthControllerTest {

    @Autowired MockMvc mockMvc;
    @Autowired ObjectMapper objectMapper;

    @Test
    void register_ShouldReturn200_AndToken() throws Exception {
        RegisterRequest req = new RegisterRequest();
        req.setUsername("user_" + System.currentTimeMillis());
        req.setEmail("e_" + System.currentTimeMillis() + "@test.pl");
        req.setPassword("haslo123");
        req.setFirstName("Jan"); req.setLastName("Nowak");

        mockMvc.perform(post("/api/auth/register")
                .contentType(MediaType.APPLICATION_JSON)
                .content(objectMapper.writeValueAsString(req)))
                .andExpect(status().isOk())
                .andExpect(jsonPath("$.token").isNotEmpty())
                .andExpect(jsonPath("$.role").value("PATIENT"));
    }
}
\end{lstlisting}

\subsection{Test frontendu -- Angular}

\begin{lstlisting}[style=java, caption={Test serwisu autoryzacji -- auth.spec.ts}]
describe('AuthService', () => {
  let service: AuthService;
  let httpMock: HttpTestingController;

  beforeEach(() => {
    TestBed.configureTestingModule({
      imports: [HttpClientTestingModule],
      providers: [AuthService]
    });
    service = TestBed.inject(AuthService);
    httpMock = TestBed.inject(HttpTestingController);
  });

  it('should store token after login', () => {
    service.login('jan', 'haslo123').subscribe();
    const req = httpMock.expectOne(
        'http://localhost:8080/api/auth/login');
    req.flush({ token: 'jwt-abc', username: 'jan', role: 'PATIENT' });

    expect(localStorage.getItem('token')).toBe('jwt-abc');
    expect(service.isLoggedIn()).toBeTrue();
  });

  afterEach(() => httpMock.verify());
});
\end{lstlisting}

% ============================================================
\section{Informacja uruchomieniowa}

\subsection{Wymagania}

\begin{itemize}
  \item Docker Engine $\geq$ 24.0 i Docker Compose $\geq$ 2.0
  \item Wolne porty: 4200 (frontend), 8080 (backend), 5432 (PostgreSQL)
\end{itemize}

\subsection{Uruchomienie przez Docker Compose}

\begin{lstlisting}[style=bash]
git clone https://github.com/REPO/Proj_ZTI.git
cd Proj_ZTI
docker compose up --build
# Frontend:  http://localhost:4200
# Backend:   http://localhost:8080
\end{lstlisting}

\begin{lstlisting}[style=bash]
docker compose down       # zatrzymuje kontenery
docker compose down -v    # zatrzymuje + usuwa wolumen bazy
\end{lstlisting}

\subsection{Zmienne środowiskowe}

\begin{table}[H]
\centering
\begin{tabularx}{\textwidth}{lXl}
\toprule
\textbf{Zmienna} & \textbf{Opis} & \textbf{Domyślna} \\
\midrule
\texttt{SPRING\_DATASOURCE\_URL}      & URL bazy PostgreSQL & jdbc:postgresql://postgres:5432/doctor\_reservation \\
\texttt{SPRING\_DATASOURCE\_USERNAME} & Użytkownik bazy     & zti\_user \\
\texttt{SPRING\_DATASOURCE\_PASSWORD} & Hasło bazy          & zti\_password \\
\texttt{JWT\_SECRET}                  & Klucz podpisu JWT   & \textbf{zmień na produkcji!} \\
\bottomrule
\end{tabularx}
\caption{Zmienne środowiskowe}
\end{table}

\subsection{Struktura projektu}

\begin{lstlisting}[style=bash]
Proj_ZTI/
|-- docker-compose.yml
|-- backend/
|   |-- Dockerfile
|   |-- pom.xml
|   +-- src/main/java/com/zti/doctorreservation/
|       |-- config/      # SecurityConfig, WebSocketConfig, DataInitializer
|       |-- model/       # User, Doctor, Patient, TimeSlot, Appointment
|       |-- repository/  # interfejsy Spring Data JPA
|       |-- service/     # AuthService, AppointmentService, TimeSlotService
|       |-- controller/  # AuthController, DoctorController, AppointmentController
|       |-- security/    # JwtUtil, JwtAuthFilter, UserDetailsServiceImpl
|       +-- aspect/      # LoggingAspect (Spring AOP)
+-- frontend/
    |-- Dockerfile
    |-- nginx.conf
    +-- src/app/
        |-- components/  # login, register, doctor-list, doctor-slots, my-appointments
        |-- services/    # auth, doctor, appointment, websocket
        |-- guards/      # authGuard
        +-- interceptors/# authInterceptor
\end{lstlisting}

% ============================================================
\section{Podręcznik użytkownika}

\subsection{Rejestracja}

\begin{enumerate}
  \item Otwórz aplikację pod adresem \texttt{http://localhost:4200}.
  \item Kliknij \textbf{Rejestracja} w pasku nawigacyjnym.
  \item Wypełnij formularz: imię, nazwisko, nazwa użytkownika, email, hasło (min. 6 znaków).
  \item Kliknij \textbf{Zarejestruj się} -- nastąpi automatyczne zalogowanie i przekierowanie do listy lekarzy.
\end{enumerate}

\subsection{Rezerwacja wizyty}

\begin{enumerate}
  \item Po zalogowaniu przejdź do \textbf{Lekarze} i wybierz lekarza.
  \item Na widoku terminów zobaczysz dostępne sloty -- statusy aktualizują się w czasie rzeczywistym:
    \begin{itemize}
      \item \textbf{Dostępny} -- można zarezerwować,
      \item \textbf{Rezerwowany\ldots} -- inny pacjent właśnie rezerwuje,
      \item \textbf{Zajęty} -- niedostępny.
    \end{itemize}
  \item Kliknij \textbf{Zarezerwuj} -- wizyta zostanie zapisana i nastąpi przekierowanie do historii wizyt.
\end{enumerate}

\subsection{Anulowanie wizyty}

\begin{enumerate}
  \item Przejdź do \textbf{Moje wizyty}.
  \item Przy zaplanowanej wizycie kliknij \textbf{Anuluj wizytę}.
  \item Slot stanie się ponownie dostępny dla innych pacjentów.
\end{enumerate}

\subsection{Dane testowe}

\begin{table}[H]
\centering
\begin{tabular}{lll}
\toprule
\textbf{Login} & \textbf{Hasło} & \textbf{Specjalizacja} \\
\midrule
dr.kowalski   & haslo123 & Kardiologia \\
dr.nowak      & haslo123 & Pediatria \\
dr.wisniewski & haslo123 & Ortopedia \\
\bottomrule
\end{tabular}
\caption{Domyślne konta lekarzy (tworzone automatycznie)}
\end{table}

% ============================================================
\section{JavaDoc -- przykłady}

\begin{lstlisting}[style=java, caption={JavaDoc -- AppointmentService.java}]
/**
 * Serwis obsługujący logikę biznesową rezerwacji wizyt.
 * Po każdej zmianie statusu slotu rozgłasza aktualizację
 * przez WebSocket do wszystkich podłączonych klientów.
 */
@Service
public class AppointmentService {

    /**
     * Rezerwuje wizytę dla zalogowanego pacjenta.
     *
     * @param request obiekt zawierający ID slotu i notatki
     * @return zapisana wizyta
     * @throws IllegalStateException gdy slot jest niedostępny
     * @throws IllegalArgumentException gdy slot nie istnieje
     */
    public Appointment book(AppointmentRequest request) { ... }

    /**
     * Anuluje wizytę i zwalnia slot czasowy.
     *
     * @param appointmentId ID wizyty do anulowania
     * @return wizyta ze statusem CANCELLED
     */
    public Appointment cancel(Long appointmentId) { ... }
}
\end{lstlisting}

\end{document}
